\section{Experimental Setup}\label{sec:setup}
\paragraph{Data} Argoverse~2 \cite{wilson2021argoverse2}, nuScenes \cite{caesar2020nuscenes} and Waymo Open Motion
\cite{ettinger2021waymo} (30 of 150 validation shards), converted with ScenarioNet \cite{li2023scenarionet} and loaded
through UniTraj \cite{feng2024unitraj}. AV2 scenarios are independent 11\,s clips with one focal agent; the labelled AV2
validation split is divided by a fixed salted hash into a calibration half and a test half. nuScenes has a mean of \nsAgentsMean\ agents per
20\,s scene (\nsNAgents\ from \nsNScenes\ validation scenes), so its calibration/test split is made by scene (\nsNScenesCal/\nsNScenesTest\ scenes). Waymo
scenarios carry \wmAgentsPerScenario\ agents on average and are split by scenario (\wmNScenariosCalAvW\ calibration, \wmNScenariosAvW\ test).
\paragraph{Forecasters} AutoBot-Ego \cite{girgis2022autobots} with $K=6$ modes. The \emph{confirmatory} models were trained
after the hypotheses about them were registered, with UniTraj's recipe (effective batch 128, learning rate
$7.5\times10^{-4}$ halved at epochs 10--50): \textbf{AB-AV2-s0} on the full 183k-scene AV2 training set and
\textbf{AB-NS-s0} on the full nuScenes training split. (Names read architecture--data--run: AB AutoBot, WF Wayformer; s0, s1 training seeds; r1--r3 reduced budget.) Each had to reach validation minADE\textsubscript{6} within
15\% of UniTraj's published value \cite{feng2024unitraj}; their final checkpoints reach \minADESixConf\ (gate 0.98, reference 0.85) and
\minADESixConfRev\ (gate 1.39, reference 1.21). The final model is the last checkpoint, fixed in advance. A second
architecture, Wayformer \cite{nayakanti2023wayformer}, was trained forward-only under the same accuracy gate and effective batch size, with its own learning rate and 27-epoch schedule
(reaching \minADESixWayformer, Sec.~\ref{sec:wf}); no reverse-direction Wayformer was trained. Four
\emph{reduced-budget} AutoBot models (training-data subsets or shorter schedules; the three AV2 models reach
minADE\textsubscript{6} \explAdeMin--\explAdeMax) were used to design the analyses and are labelled exploratory throughout.
\paragraph{Protocol} $\alpha=\delta=0.1$. Forward: calibrate on the AV2 calibration half; test on the AV2 test half
(same-domain control), nuScenes and Waymo. Reverse: calibrate on the nuScenes calibration scenes; test on nuScenes test
scenes, AV2 and Waymo. Label-budget experiments draw $k$ agents or $m$ whole scenes 1000 times from a labelled pool that
is disjoint from the evaluation pool. Intervals on multi-agent targets resample scenes. The source threshold is split conformal on all calibration scores: for AV2 each
independent clip contributes one score, so its guarantee applies as stated; for nuScenes and Waymo the scores of agents in one
scene are dependent, so the formal guarantee does not cover the source threshold itself, and we verify it empirically on a
held-out same-domain split ($\gapSameConf$ and $\gapSameConfRev$ points) and measure every loss against that control.
\paragraph{Registration and disclosures} Hypotheses, thresholds and success criteria were written to a version-controlled,
timestamped log (released with the code) before the runs they govern; later analyses are dated addenda and outcomes are
appended, never edited. Disclosed defects: a cache collision, per-agent nuScenes splits and truncated scenario
identifiers (none affecting a reported number), and a betting $p$-value whose bet cap allowed negative capital, found in
an internal review, covered by a regression test, and followed by recomputing every repair result in
Sec.~\ref{sec:repair-scene} under a design registered beforehand.
